\section{Conclusão}

\paragraph{}
Ao longo desta atividade foi possível compreender a importância da etapa de
enumeração durante processos de análise de segurança e testes de invasão. A
utilização de ferramentas automatizadas demonstrou como grandes volumes de
informações podem ser coletados rapidamente, permitindo identificar potenciais
vetores de escalação de privilégios de forma mais eficiente.

\paragraph{}
As práticas realizadas evidenciaram diferentes categorias de vulnerabilidades
comumente encontradas em ambientes Windows, incluindo permissões excessivas em
serviços, configurações inadequadas de caminhos de execução, carregamento
inseguro de bibliotecas dinâmicas e exposição de arquivos críticos contendo
informações de autenticação.

\paragraph{}
Também foi possível observar como pequenas falhas administrativas podem
resultar na execução de comandos com privilégios elevados, possibilitando a
concessão de acesso administrativo a usuários inicialmente restritos.

\paragraph{}
Por fim, a atividade reforçou a necessidade da aplicação de boas práticas de
segurança, revisões periódicas de configuração e monitoramento contínuo dos
sistemas, reduzindo a superfície de ataque e dificultando a exploração dessas
vulnerabilidades por agentes maliciosos.
